\section{Discussion and Limitations}
\label{sec:disc}
\paragraph{What the results say.}
On both leaderboards, first place under the deployed averaging rules is cheap in tasks for top-ranked targets and expensive in points: a single sub-dataset can suffice, but the points needed under the mean are $n$ times the gap to the leader and grow steeply with the target's rank (log--log slope 0.4--0.7 on OLLB and 1.1--2.2 on HELM at \texttt{ceil}, as a descriptive summary). Whether other rules are easier or harder depends on what one counts. Median, minimum and plurality need far fewer points, and targeting pays for them (cost 3--51\% of the uniform improvement), while the mean is indifferent to where the points go; rules that need a majority of pairwise wins (Borda, Condorcet requirement) need many tasks. Realistic caps change the picture more than the rule does: under \texttt{emp95} or \texttt{gain10}, at most 2.5\% of the OLLB targets beyond rank 500 can win under any score rule (none under the mean), but \texttt{gain20} lets 60\% of them win under the minimum rule. The results suggest that a designer who wants a leaderboard to resist targeted investment would do better to bound what can be gained per sub-dataset (for example by holding data out) than to rely on the choice of aggregation rule alone; we did not test such a defence beyond the clipping of Appendix~\ref{app:checks}, which leaves the cost of first place unchanged, and we do not claim that any rule is robust.

\paragraph{Limitations.}
(i) \emph{The cap is an assumption.} The empirical cap uses 540 fine-tune/base pairs of one leaderboard, 154 labelled chat models and 346 fine-tuned on domain-specific datasets, so it mixes generic instruction tuning with the targeted training we want to model (dropping chat-typed pairs changes the feasible shares by under a point, but they range from 0.4 to 16\% with the rule and the definition, and the minimum rule reaches 100\% under a flat quarter of the headroom; Appendix~\ref{app:checks}); it is the same for every target whatever its starting score; and it says nothing about what a determined developer with new data could reach (the leakage studies we cite report larger gains), so it may understate as well as overstate. ``Investment'' is score points, not data or compute. (ii) \emph{Single deviator.} Other models do not react, and several developers acting together or a designer changing the rule is outside the model. (iii) \emph{Small $n$ on OLLB.} With six tasks, rank rules are coarse, and the even-$n$ median has no simple closed form; the HELM results (one benchmark family, 79 models, 57 correlated subjects with 100 to 1534 test items each, so one item is 1.0 point at the smallest) should not be generalised to other leaderboards, and $K$ counts six distinct benchmarks on OLLB but 57 subjects of one benchmark on HELM, so we compare $K$ within a dataset, not across. (iv) \emph{Sampled ordinal rules.} Borda, Condorcet requirement and plurality on OLLB were solved for a stratified sample of 136 targets, not all 4496, and no exact Copeland, Schulze or Kemeny ranking was computed; the Condorcet requirement is only an upper bound for those rules. (v) \emph{Stability.} Rule orders differ between the two leaderboards and, on OLLB with six tasks, under task resampling, and agree only moderately across caps (Kendall $\tau$ about 0.5; Section~\ref{sec:results}); the $\tau$ values have no null distribution (not computed), and individual targets can change cost substantially with the tie margin. (vi) \emph{Scores are frozen snapshots}, the data carry no declared licence, and the models' scores were taken as given, including possible contamination. (vii) \emph{Compute:} 7.7 CPU-hours in total, above the task's rough 4-hour estimate (Appendix~\ref{app:repro}).

\paragraph{Broader impact.}
The analysis describes how cheaply rankings can be moved by targeted training, which could inform gaming; it uses only public scores and trains no model, and its practical content is for benchmark designers (cap what a sub-dataset can contribute, report intervals). We make no claim about any real developer.
